\chapter{Groups, Rings, and Fields}

By now you have spent two chapters (Chapters~3 and~4) living inside the integers modulo $n$: adding, multiplying, finding inverses, hunting for primes. You may have noticed that the same moves keep coming back. Addition mod $n$ behaves a lot like ordinary addition. The nonzero residues mod a prime behave a lot like ordinary multiplication. Inverses appear exactly when a $\gcd$ is $1$. These are not coincidences. They are symptoms of three \emph{algebraic structures} that crypto is built on top of: groups, rings, and fields.

This chapter is the bridge between ``I can compute mod $n$'' and ``I understand why
Diffie--Hellman, RSA, and elliptic curves work.'' We are going to name the structures,
state their rules precisely, and -- this is the part most courses rush -- build the
intuition for \emph{why} a cryptographer cares. The single most important payoff is
Lagrange's theorem, which tells you that the order of an element must divide the size of
the group. That one fact is the engine behind RSA's correctness, the security parameters
of Diffie--Hellman, and the choice of curve parameters in ECC. We will get there
carefully, with tiny numbers you can check on your fingers.

\chapterline

\section{Why crypto needs algebra at all}
\label{ch05:why-algebra}

Cryptography is the art of doing arithmetic that is \emph{easy one way and hard the
other}. You want a function the legitimate parties can evaluate quickly and the
adversary cannot invert. Plain integer arithmetic is hopeless for this: addition,
multiplication, and even exponentiation over $\Z$ are all easy to undo. There is nothing
to hide behind.

The trick is to do arithmetic inside a \emph{finite} world that wraps around. When you
work modulo $n$, the operation $g \mapsto g^x \bmod n$ scrambles the order of elements so
thoroughly that recovering $x$ from $g^x$ -- the discrete logarithm problem -- has no
known fast algorithm for well-chosen $n$. The wrapping-around is what creates the
asymmetry between ``forward'' (cheap) and ``backward'' (apparently astronomically
expensive).

But the moment you commit to a finite, wrap-around arithmetic, you need to know which of
your familiar rules survive. Can you always divide? When does an inverse exist? How many
times can you square an element before it cycles back to where it started? Those
questions have clean answers, and the answers are exactly the theorems about groups,
rings, and fields. So algebra is not decoration bolted onto cryptography. It is the
vocabulary in which ``this is secure'' becomes a statement you can actually prove.

\begin{intuitionbox}
Think of a structure as a \emph{game with rules}. A group is the minimal game where you
can combine two things, undo any move, and never fall off the board. A ring adds a second
operation so you have both ``add'' and ``multiply.'' A field is the luxury edition where
you can also \emph{divide} by anything nonzero. Crypto schemes are recipes written in
this game; the security proofs are arguments about what the rules do and do not let an
attacker do.
\end{intuitionbox}

\section{Groups}
\label{ch05:groups}

\subsection{The definition, and why each axiom earns its place}

\begin{definition}[Group]
\label{ch05:def-group}
A \emph{group} is a set $G$ together with a binary operation $\ast : G \times G \to G$
satisfying four axioms:
\begin{enumerate}
  \item \textbf{Closure.} For all $a,b \in G$, the result $a \ast b$ is again in $G$.
  \item \textbf{Associativity.} For all $a,b,c \in G$, $(a \ast b)\ast c = a \ast (b \ast c)$.
  \item \textbf{Identity.} There is an element $e \in G$ with $e \ast a = a \ast e = a$ for every $a$.
  \item \textbf{Inverses.} For every $a \in G$ there is some $a^{-1} \in G$ with
        $a \ast a^{-1} = a^{-1} \ast a = e$.
\end{enumerate}
If in addition $a \ast b = b \ast a$ for all $a,b$, the group is \emph{abelian}
(commutative). Every group in this book is finite and abelian unless we say otherwise.
\end{definition}

Each axiom pulls its weight. Closure says the operation never throws you out of the
set -- you can compose moves forever and stay inside the game. Associativity is what lets
you write $a \ast b \ast c \ast d$ without parentheses, which is what makes
\emph{repeated} operation (and hence exponentiation) well defined. Identity gives you a
``do nothing'' move, the anchor that ``$g^0$'' or ``$0 \cdot P$'' refers to. Inverses are
the whole reason groups matter to crypto: every action is reversible, so encryption with a
group operation can always be undone by the right party.

\begin{focusbox}
Notice what is \emph{not} required: there is no second operation, no notion of ``bigger''
or ``smaller,'' no division. A group is deliberately spare. That spareness is a feature.
Because the axioms are so few, an enormous range of objects qualify -- integers under
addition, nonzero residues under multiplication, points on an elliptic curve, symmetries
of a cube -- and any theorem proved from the four axioms applies to all of them at once.
\end{focusbox}

\subsection{The two groups you already know}

You have met two groups already without the name.

\textbf{Additive group $\Zn{n}$.} Take the residues $\{0,1,\dots,n-1\}$ with the
operation ``add, then reduce mod $n$.'' The identity is $0$; the inverse of $a$ is $n-a$
(and $0$ is its own inverse). Closure and associativity come straight from ordinary
addition. This group always has exactly $n$ elements, for every $n \ge 1$.

\textbf{Multiplicative group $\Zns{n}$.} Take the residues that are \emph{coprime} to $n$,
with operation ``multiply, then reduce mod $n$.'' The identity is $1$. The catch -- and
the reason we throw away the non-coprime residues -- is inverses. Recall from the previous
chapter that $a$ has a multiplicative inverse mod $n$ if and only if $\gcd(a,n)=1$. So if
we want \emph{every} element to be invertible (axiom 4), we must keep only the coprime
ones. The size of this group is Euler's totient $\varphi(n)$.

\begin{example}[$\Zns{15}$ by hand]
\label{ch05:ex-zns15}
The residues coprime to $15 = 3\cdot 5$ are
$\{1,2,4,7,8,11,13,14\}$, so $\varphi(15)=8$. Check closure on a couple of products:
$2 \cdot 8 = 16 \equiv 1$, $4 \cdot 4 = 16 \equiv 1$, $7 \cdot 13 = 91 = 6\cdot 15 + 1
\equiv 1$. Each of those pairs is mutually inverse. Notice $5 \notin \Zns{15}$ because
$\gcd(5,15)=5 \neq 1$: indeed $5 \cdot 3 = 15 \equiv 0$, and an element that can multiply
to $0$ can never have an inverse.
\end{example}

\begin{trapbox}
A common slip: thinking $\Zns{n}$ is ``all of $\{1,\dots,n-1\}$.'' That is only true when
$n$ is prime. For composite $n$ you must drop every residue sharing a factor with $n$.
Forgetting this is the source of countless ``why is my inverse code returning garbage''
bugs -- the modular inverse simply does not exist for those inputs.
\end{trapbox}

\subsection{Group notation: additive vs. multiplicative}

The same abstract group can be written two ways, and crypto uses both. When we write the
operation as ``$+$'', repeating it $k$ times is written $k \cdot a = a + a + \cdots + a$,
the identity is $0$, and the inverse is $-a$. When we write the operation as
``$\cdot$'' (or just juxtaposition), repeating it $k$ times is written $a^k = a \cdot a
\cdots a$, the identity is $1$, and the inverse is $a^{-1}$. Elliptic-curve people use
additive notation (so the secret-key operation is ``multiply a point by a scalar''); the
$\Zns{p}$ world uses multiplicative notation (so the operation is ``raise to a power'').
Both notations describe the exact same group operations.

\begin{walkbox}
Quick check before moving on: in the additive group $\Zn{6}$, what is the inverse of $4$?
You need $x$ with $4 + x \equiv 0 \pmod 6$, so $x = 2$. In the multiplicative group
$\Zns{7}$, what is the inverse of $4$? You need $4x \equiv 1 \pmod 7$; testing,
$4 \cdot 2 = 8 \equiv 1$, so the inverse is $2$. Same word ``inverse,'' two different
operations.
\end{walkbox}

\section{Order of an element, and cyclic groups}
\label{ch05:order}

\subsection{Order of an element}

Pick a group element $g$ and start applying the operation to itself: $g, g^2, g^3,
\dots$ (in multiplicative notation). Because the group is finite, this list cannot go on
forever without repeating. The first time it lands back on the identity is the element's
\emph{order}.

\begin{definition}[Order of an element]
\label{ch05:def-order}
The \emph{order} of $g \in G$, written $\ord(g)$, is the smallest positive integer $k$
such that $g^k = e$. (In additive notation: the smallest $k>0$ with $k\cdot g = 0$.)
\end{definition}

\begin{example}[Orders in $\Zns{7}$]
\label{ch05:ex-orders-7}
Here $7$ is prime so $\Zns{7} = \{1,2,3,4,5,6\}$, six elements. Let us trace the powers of
a few:
\[
\begin{array}{l|cccccc}
g & g^1 & g^2 & g^3 & g^4 & g^5 & g^6 \\\hline
2 & 2 & 4 & 1 & 2 & 4 & 1 \\
3 & 3 & 2 & 6 & 4 & 5 & 1 \\
6 & 6 & 1 & 6 & 1 & 6 & 1
\end{array}
\]
Reading off the first return to $1$: $\ord(2)=3$, $\ord(3)=6$, $\ord(6)=2$. Also
$\ord(1)=1$ since $1^1 = 1$. Notice every order we found -- $1,2,3,6$ -- divides $6$, the size
of the group. This illustrates Lagrange's theorem in miniature: an element's orbit length must cleanly divide the total order of the group.
\end{example}

\begin{intuitionbox}
The powers of $g$ form a \emph{cycle}: a closed loop that returns to the identity and then
repeats. The order is the length of that loop. A short loop means $g$ ``can't reach much''
by exponentiation; a long loop means $g$ sweeps through many elements before coming home.
For crypto we usually want long loops -- an element whose powers cover as much ground as
possible, so that $g^x$ is a good hiding place for the secret $x$.
\end{intuitionbox}

\subsection{Cyclic groups and generators}

When a single element's loop covers the \emph{entire} group, that element is a generator
and the group is cyclic.

\begin{definition}[Cyclic group, generator]
\label{ch05:def-cyclic}
A group $G$ is \emph{cyclic} if there exists $g \in G$ -- a \emph{generator} -- such that
every element of $G$ equals $g^k$ for some integer $k$. Equivalently, $\ord(g) = \abs{G}$.
We write $G = \langle g \rangle$.
\end{definition}

In Example~\ref{ch05:ex-orders-7}, the element $3$ has order $6 = \abs{\Zns{7}}$, so its
powers $3,2,6,4,5,1$ run through all six elements. Thus $\Zns{7}$ is cyclic with generator
$3$ (Figure~\ref{ch05:fig-cycle}). The element $2$, by contrast, has order $3$ and
generates only the smaller set $\{1,2,4\}$ -- a perfectly good group in its own right,
sitting inside the big one.

\begin{figure}[H]
\centering
\begin{tikzpicture}[font=\small,>=stealth,scale=1.05]
  \def\R{2.2}

  % 6 nodes positioned around the circle
  % Angle = 90 - i * 60:
  % i=0 (90 deg): 1
  % i=1 (30 deg): 3
  % i=2 (-30 deg): 2
  % i=3 (-90 deg): 6
  % i=4 (-150 deg): 4
  % i=5 (150 deg): 5
  \foreach \ang/\val/\col in {90/1/structurepurple, 30/3/sourceblue, -30/2/targetgreen, -90/6/quotientorange, -150/4/targetgreen, 150/5/sourceblue}{
    \node[circle,draw=\col!80,thick,fill=\col!12,minimum size=8mm,font=\bfseries]
      (n\val) at (\ang:\R) {$\val$};
  }

  % Outer master cycle: <3> of order 6 (Blue arrows, x3)
  \foreach \a/\b in {1/3, 3/2, 2/6, 6/4, 4/5, 5/1}{
    \draw[->,very thick,sourceblue!80] (n\a) to[bend left=20] (n\b);
  }

  % Order-3 subgroup <2> = {1, 2, 4} (Green chords, x2)
  \draw[->,thick,dashed,targetgreen!90!black] (n1) to[bend right=10] (n2);
  \draw[->,thick,dashed,targetgreen!90!black] (n2) to[bend right=10] (n4);
  \draw[->,thick,dashed,targetgreen!90!black] (n4) to[bend right=10] (n1);

  % Order-2 subgroup <6> = {1, 6} (Orange straight diameter, x6)
  \draw[<->,thick,densely dotted,quotientorange!90!black] (n1) -- (n6);

  % Center annotations
  \node[font=\footnotesize\bfseries,align=center] at (0.2,0.3) {Master Orbit\\$\langle 3 \rangle = \Zns{7}$};
  \node[font=\tiny,gray] at (0.2,-0.2) {$\abs{\Zns{7}} = 6$};

  % Legend below
  \node[draw=gray!40,fill=gray!5,rounded corners=4pt,thick,align=left,font=\tiny,inner sep=5pt] at (0,-3.3) {
    \textbf{Subgroup Decomposition (Lagrange in Action):}\\
    $\bullet$ \textcolor{sourceblue!90!black}{\textbf{Solid Blue Outer Ring:}} Generator $\langle 3 \rangle$ sweeps full order 6 ($1\to 3\to 2\to 6\to 4\to 5\to 1$).\\
    $\bullet$ \textcolor{targetgreen!90!black}{\textbf{Dashed Green Triangle:}} Subgroup $\langle 2 \rangle = \{1, 2, 4\}$ of order 3 ($3 \mid 6$).\\
    $\bullet$ \textcolor{quotientorange!90!black}{\textbf{Dotted Orange Diameter:}} Subgroup $\langle 6 \rangle = \{1, 6\}$ of order 2 ($2 \mid 6$).\\
    $\bullet$ Identity element $1$ anchors every subgroup orbit.
  };
\end{tikzpicture}
\caption{The cyclic group $\Zns{7} = \langle 3 \rangle$ and its internal subgroup orbits. The generator 3 visits all six elements. Inside this orbit live the order-3 subgroup $\langle 2 \rangle = \{1,2,4\}$ and the order-2 subgroup $\langle 6 \rangle = \{1,6\}$. Every subgroup orbit length divides the master group size $6 = \varphi(7)$.}
\label{ch05:fig-cycle}
\end{figure}

\begin{example}[$\Zn{6}$ is cyclic too]
\label{ch05:ex-z6-cyclic}
In additive notation, ``powers of $g$'' means multiples. Take $g=1$ in $\Zn{6}$:
$1,2,3,4,5,0$ -- all six residues, so $1$ generates. Take $g=5$: $5,4,3,2,1,0$ -- also all
six, so $5$ generates too. But $g=2$ gives $2,4,0,2,4,0,\dots$, only $\{0,2,4\}$: order
$3$, not a generator. The generators of $\Zn{6}$ are exactly $1$ and $5$, the residues
coprime to $6$. (That is a general fact: $k$ generates $\Zn{n}$ iff $\gcd(k,n)=1$.)
\end{example}

\begin{whybox}
Why cyclic groups are the crypto sweet spot. In a cyclic group $\langle g \rangle$ of size
$q$, every element is $g^x$ for a unique $x \in \{0,1,\dots,q-1\}$. The map $x \mapsto g^x$
is therefore a bijection between exponents and group elements. Going \emph{forward}
($x \mapsto g^x$) is a handful of multiplications via fast exponentiation. Going
\emph{backward} ($g^x \mapsto x$, the discrete logarithm) has no known efficient
algorithm in well-chosen groups. That gap -- a bijection that is easy one way and hard the
other -- is precisely a one-way function with structure, and structure is what lets two
parties combine their secrets (Diffie--Hellman) without ever revealing them.
\end{whybox}

\begin{historybox}
The word \emph{group} comes from \'Evariste Galois, who in 1832 -- the night before a fatal duel at age twenty -- studied permutations of polynomial roots to characterize which equations can be solved by radicals; rejected by the French Academy during his lifetime, his manuscripts were published by Joseph Liouville in 1846 and founded abstract algebra. A century and a half later, Whitfield Diffie and Martin Hellman's 1976 paper ``New Directions in Cryptography'' put cyclic groups at the heart of public-key cryptography. They worked in $\Zns{p}$ for a large prime $p$. Decades of effort have made the discrete log \emph{somewhat} easier in $\Zns{p}$ (index-calculus methods), which is why modern systems migrated to elliptic-curve groups, where those methods do not apply and the same security needs far smaller numbers. The underlying algebra -- a cyclic group with a hard discrete log -- is identical; only the group changed.
\end{historybox}

\subsection{Not every group is cyclic}

$\Zns{n}$ is cyclic for some $n$ and not for others. It \emph{is} cyclic exactly when
$n \in \{1,2,4,p^k,2p^k\}$ for an odd prime $p$. The smallest cautionary example is
$\Zns{8} = \{1,3,5,7\}$.

\begin{example}[$\Zns{8}$ has no generator]
\label{ch05:ex-zns8}
Square every element mod $8$: $1^2=1$, $3^2=9\equiv 1$, $5^2=25\equiv 1$, $7^2=49\equiv
1$. \emph{Every} element squares to $1$, so every non-identity element has order $2$. No
element has order $4 = \abs{\Zns{8}}$, hence no generator exists. This group is not
cyclic; structurally it is the ``Klein four-group,'' two independent order-$2$ switches.
\end{example}

This is not a curiosity. When a protocol assumes a generator exists and you hand it a
non-cyclic group, the math quietly breaks. Picking the modulus so that $\Zns{n}$ is cyclic
(or, more precisely, working inside a prime-order cyclic \emph{subgroup}) is part of doing
crypto correctly, which we make precise after Lagrange.

\section{Subgroups and Lagrange's theorem}
\label{ch05:lagrange}

\subsection{Subgroups}

A \emph{subgroup} is a subset that is itself a group under the same operation.

\begin{definition}[Subgroup]
\label{ch05:def-subgroup}
A nonempty subset $H \subseteq G$ is a \emph{subgroup} (written $H \le G$) if it is closed
under the operation and under inverses: $a,b \in H \Rightarrow a \ast b \in H$ and
$a^{-1} \in H$. (These two conditions force $e \in H$ and inherit associativity from $G$.)
\end{definition}

The cleanest source of subgroups is a single element. The set of all powers of $g$,
\[
\langle g \rangle = \{e, g, g^2, \dots, g^{\ord(g)-1}\},
\]
is always a subgroup -- the \emph{cyclic subgroup generated by $g$}. Its size is exactly
$\ord(g)$. We saw this concretely: in $\Zns{7}$, $\langle 2 \rangle = \{1,2,4\}$ has
$3 = \ord(2)$ elements.

\subsection{Lagrange's theorem}

Here is the structural fact that everything else in this chapter has been building toward.

\begin{theorem}[Lagrange]
\label{ch05:thm-lagrange}
If $G$ is a finite group and $H \le G$ is a subgroup, then $\abs{H}$ divides $\abs{G}$.
\end{theorem}

\begin{whybox}
Why Lagrange is true (the coset argument). For a fixed $a \in G$, define the
\emph{coset} $aH = \{a \ast h : h \in H\}$. Two facts do all the work.

\emph{Every coset has the same size as $H$.} The map $h \mapsto a \ast h$ from $H$ to
$aH$ is reversible -- multiply by $a^{-1}$ to undo it -- so it is a bijection. Hence
$\abs{aH} = \abs{H}$.

\emph{Distinct cosets do not overlap.} Suppose $aH$ and $bH$ share an element. Then
$a h_1 = b h_2$ for some $h_1, h_2 \in H$, so $b^{-1}a = h_2 h_1^{-1} \in H$. Multiplying
$H$ on the left by an element of $H$ just reshuffles $H$, which forces $aH = bH$. So any
two cosets are either identical or disjoint.

Now every element $a$ lives in its own coset $aH$ (since $e \in H$ gives
$a = a\ast e \in aH$), so the cosets \emph{tile} $G$ completely: they partition it into
disjoint blocks, each of size $\abs{H}$. If there are $m$ blocks, then
$\abs{G} = m\cdot\abs{H}$, and therefore $\abs{H}$ divides $\abs{G}$. The number $m$ is
called the \emph{index} $[G:H]$.
\end{whybox}

\begin{figure}[H]
\centering
\begin{tikzpicture}[font=\small,>=stealth,scale=0.95]
  \def\bw{2.6}

  % Coset Cards
  \foreach \i/\lab/\col/\rep in {
    0/{$eH = H$}/sourceblue/e,
    1/{$aH$}/targetgreen/a,
    2/{$bH$}/quotientorange/b,
    3/{$cH$}/structurepurple/c
  } {
    \draw[\col!70, thick, fill=\col!10, rounded corners=4pt] (\i*\bw, 0) rectangle (\i*\bw+\bw-0.35, 2.4);
    \node[\col!90!black, font=\bfseries\footnotesize] at (\i*\bw+\bw/2-0.17, 2.05) {\lab};

    % Discrete elements
    \node[draw=\col!80, fill=white, rounded corners=2pt, font=\tiny, align=center, inner sep=2pt] at (\i*\bw+\bw/2-0.17, 1.25) {
      $\rep \ast h_1, \;\rep \ast h_2$\\
      $\rep \ast h_3, \;\rep \ast h_4$
    };

    \node[\col!80!black, font=\tiny] at (\i*\bw+\bw/2-0.17, 0.4) {size $= |H|$};
  }

  % Bijection arrows between first coset and others
  \draw[->,thick,dashed,targetgreen!80] (1.12, 2.48) to[bend left=22] node[pos=0.52, fill=white, inner sep=1pt, font=\tiny, text=targetgreen!90!black]{$\times a$} (3.72, 2.48);
  \draw[->,thick,dashed,quotientorange!80] (1.12, 2.52) to[bend left=32] node[pos=0.52, fill=white, inner sep=1pt, font=\tiny, text=quotientorange!90!black]{$\times b$} (6.32, 2.52);
  \draw[->,thick,dashed,structurepurple!80] (1.12, 2.56) to[bend left=40] node[pos=0.52, fill=white, inner sep=1pt, font=\tiny, text=structurepurple!90!black]{$\times c$} (8.92, 2.56);

  % Bottom brace
  \draw[decorate,decoration={brace,amplitude=8pt,mirror},thick,black] (-0.1,-0.3) -- (4*\bw-0.45,-0.3)
    node[midway,below=10pt,font=\footnotesize\bfseries,align=center] {
      Partition of Group $G$ into $m = [G:H]$ Disjoint Cosets\\
      $|G| = m \cdot |H| \implies |H| \text{ cleanly divides } |G|$
    };
\end{tikzpicture}
\caption{Visual proof of Lagrange's theorem by coset tiling. The subgroup $H$ is mapped bijectively to each disjoint coset $aH$ by left multiplication. Because every element of $G$ belongs to exactly one coset and every coset has identical size $|H|$, the subgroup size $|H|$ must divide the group size $|G|$.}
\label{ch05:fig-cosets}
\end{figure}

\subsection{The consequences crypto actually uses}

Lagrange by itself is a counting statement. Its corollaries are what you reach for daily.

\begin{corollary}[Order divides group size]
\label{ch05:cor-order-divides}
For any $g$ in a finite group $G$, $\ord(g)$ divides $\abs{G}$.
\end{corollary}
This is immediate: $\ord(g) = \abs{\langle g\rangle}$ and $\langle g\rangle \le G$, so
apply Lagrange. It explains the pattern from Example~\ref{ch05:ex-orders-7}: in a group of
$6$ elements, the only possible orders are $1,2,3,6$ -- the divisors of $6$. You will never
find an element of order $4$ or $5$ there, no matter how hard you look.

\begin{corollary}[Identity at the group size]
\label{ch05:cor-glag}
For any $g$ in a finite group $G$ of order $n=\abs{G}$, we have $g^{n} = e$.
\end{corollary}
Write $k = \ord(g)$, which divides $n$ by the previous corollary, say $n = k m$. Then
$g^{n} = (g^{k})^{m} = e^{m} = e$. Raising \emph{any} element to the size of the group
sends it home to the identity.

This last corollary is a powerhouse. Specialize $G = \Zns{n}$, whose size is
$\varphi(n)$, and Corollary~\ref{ch05:cor-glag} becomes Euler's theorem; specialize
further to a prime modulus and it becomes Fermat's little theorem.

\begin{corollary}[Euler and Fermat, for free]
\label{ch05:cor-euler-fermat}
If $\gcd(a,n)=1$ then $a^{\varphi(n)} \equiv 1 \pmod n$ \textup{(Euler)}. If $p$ is prime
and $p \nmid a$ then $a^{p-1} \equiv 1 \pmod p$ \textup{(Fermat)}.
\end{corollary}

\begin{intuitionbox}
You may have seen Euler's and Fermat's theorems proved by clever direct manipulation in
the previous chapter. Here they fell out of a single structural idea -- ``order divides
group size'' -- with no manipulation at all. That is the recurring profit of the algebraic
viewpoint: a theorem that looked like a number-theory trick is revealed as an instance of
something that holds in \emph{every} finite group, elliptic curves included.
\end{intuitionbox}

\begin{example}[Euler in action]
\label{ch05:ex-euler}
Take $n=15$, where $\varphi(15)=8$ (Example~\ref{ch05:ex-zns15}). Euler promises
$a^{8}\equiv 1$ for any $a$ coprime to $15$. Check $a=2$: $2^4 = 16 \equiv 1$, and indeed
$2^8 = (2^4)^2 \equiv 1$. Here $\ord(2)=4$, which divides $8$ -- Lagrange again. So you
do not even need the full exponent $8$; the true order $4$ already does the job, but $8$
is guaranteed to work for \emph{every} element at once.
\end{example}

\begin{connectionbox}
The cost of computing $g^x \bmod n$ is the subject of the companion algorithms book's
chapter on \emph{modular exponentiation}: square-and-multiply does it in about $\log_2 x$
modular multiplications, so even a $2048$-bit exponent costs only a few thousand
multiplies. Lagrange tells you something orthogonal and equally practical: you may first
reduce the exponent modulo $\ord(g)$ (or modulo $\varphi(n)$) before you start, because
$g^{x} = g^{x \bmod \ord(g)}$. The two facts together are why RSA decryption can use a
private exponent reduced mod $\varphi(n)$ rather than something astronomically large.
\end{connectionbox}

\section{From one operation to two: rings}
\label{ch05:rings}

A group gives you one operation. Ordinary arithmetic has two -- you add \emph{and}
multiply -- and they interact through distributivity. The structure that captures ``two
compatible operations'' is a ring.

\begin{definition}[Commutative ring with identity]
\label{ch05:def-ring}
A \emph{ring} is a set $R$ with two operations $+$ and $\cdot$ such that:
\begin{itemize}
  \item $(R,+)$ is an abelian group, with additive identity $0$ and negatives $-a$;
  \item multiplication is associative, commutative, and has an identity $1$;
  \item the two are linked by \emph{distributivity}: $a\cdot(b+c) = a\cdot b + a\cdot c$.
\end{itemize}
\end{definition}

The flagship example is $\Zn{n}$ \emph{with both operations}: add mod $n$, multiply mod
$n$. Every $\Zn{n}$ is a ring. So is the set of integers $\Z$, and the set of polynomials
with coefficients in any ring -- a fact we rely on in Chapter~6,
where we build $\F_{2^8}$ (the byte field inside AES) out of polynomials mod $2$.

\begin{focusbox}
What a ring \emph{lacks} is the interesting part. Multiplication is \emph{not} required to
have inverses. In $\Zn{n}$ you can always add, subtract, and multiply, but you cannot
always divide. The elements you \emph{can} divide by are exactly the \emph{units}: those
with a multiplicative inverse, i.e. the coprime residues. Collect the units of a ring and
they form a group under multiplication -- and for $\Zn{n}$ that group is precisely
$\Zns{n}$ from Section~\ref{ch05:groups}. The ring contains a hidden group of invertible
elements.
\end{focusbox}

\subsection{Zero divisors: the wrinkle that composite moduli introduce}

In ordinary integers, if $a\cdot b = 0$ then $a=0$ or $b=0$. In a general ring that can
fail, and the failure is exactly what makes composite moduli awkward.

\begin{definition}[Zero divisor]
\label{ch05:def-zerodiv}
A nonzero $a \in R$ is a \emph{zero divisor} if there is a nonzero $b$ with $a\cdot b = 0$.
\end{definition}

\begin{example}[Zero divisors mod $12$]
\label{ch05:ex-zerodiv-12}
In $\Zn{12}$, $3\cdot 4 = 12 \equiv 0$, yet neither $3$ nor $4$ is zero. Both are zero
divisors. They are also exactly the residues sharing a factor with $12$ -- the
non-units. A zero divisor can never be a unit: if $a b = 0$ with $b \ne 0$ and $a$ also
had an inverse $a^{-1}$, then $b = a^{-1}(ab) = a^{-1}\cdot 0 = 0$, a contradiction. So in
$\Zn{n}$ each nonzero residue is \emph{either} a unit \emph{or} a zero divisor, decided by
$\gcd$ with $n$.
\end{example}

\begin{trapbox}
Zero divisors quietly wreck the cancellation law. Over the rationals, $a x = a y$ with
$a \ne 0$ lets you cancel to $x=y$. Over $\Zn{12}$ it does not: $3\cdot 2 = 6 = 3\cdot 6$,
but $2 \ne 6$. Any algorithm that ``divides through'' by a value must first check that the
value is a unit. This is a real, recurring implementation bug -- and it is the structural
reason cryptographers prefer prime moduli, where the next section shows there are no zero
divisors at all.
\end{trapbox}

\section{The luxury edition: fields}
\label{ch05:fields}

A field is a ring in which division by anything nonzero is allowed. It is the structure
where all four arithmetic operations behave the way your school algebra led you to expect.

\begin{definition}[Field]
\label{ch05:def-field}
A \emph{field} is a commutative ring with $1 \ne 0$ in which \emph{every} nonzero element
has a multiplicative inverse. Equivalently: $(F,+)$ is an abelian group and
$(F\setminus\{0\},\cdot)$ is also an abelian group.
\end{definition}

The familiar infinite fields are the rationals $\Q$, the reals $\R$, and the complex
numbers $\C$. But crypto runs on \emph{finite} fields, and the most important family is
straightforward to describe.

\begin{theorem}[Prime fields]
\label{ch05:thm-prime-field}
$\Zn{p}$ is a field if and only if $p$ is prime. When $p$ is prime we write it $\F_p$.
\end{theorem}

\begin{whybox}
Why primality is exactly the dividing line. If $p$ is prime, then every residue
$a \in \{1,\dots,p-1\}$ satisfies $\gcd(a,p)=1$, so every nonzero element is a unit and has
an inverse -- that is the definition of a field. If $p$ is composite, write $p = ab$ with
$1 < a,b < p$; then $a$ and $b$ are zero divisors (Example~\ref{ch05:ex-zerodiv-12}
showed the pattern), and a zero divisor cannot have an inverse, so $\Zn{p}$ falls short of
being a field. Primality is not a convenience here; it is the precise condition under
which ``divide by anything nonzero'' becomes legal.
\end{whybox}

\begin{example}[$\F_5$ in full]
\label{ch05:ex-f5}
The field $\F_5 = \{0,1,2,3,4\}$. Its multiplication table on nonzero elements reveals why division works:
\begin{figure}[H]
\centering
\begin{tikzpicture}[font=\small,>=stealth]
  % Define the grid coordinates
  \def\w{0.8}
  \def\h{0.7}

  % Draw grid lines
  \draw[thick] (0, 0) -- (0, -5*\h);
  \draw[thick] (0, -\h) -- (5*\w, -\h);
  
  % Header Row & Col
  \node[font=\bfseries] at (\w/2, -\h/2) {$\cdot$};
  \foreach \i in {1,2,3,4} {
    \node[font=\bfseries,sourceblue!90!black] (c\i) at (\i*\w + \w/2, -\h/2) {\i};
    \node[font=\bfseries,sourceblue!90!black] (r\i) at (\w/2, -\i*\h - \h/2) {\i};
  }

  % Inner grid (data)
  \foreach \r/\vA/\vB/\vC/\vD in {1/1/2/3/4, 2/2/4/1/3, 3/3/1/4/2, 4/4/3/2/1} {
    \foreach \c/\val in {1/\vA, 2/\vB, 3/\vC, 4/\vD} {
      % Highlight the 1s
      \ifnum\val=1
        \fill[targetgreen!20,rounded corners=2pt] (\c*\w+0.05, -\r*\h-\h+0.05) rectangle (\c*\w+\w-0.05, -\r*\h-0.05);
        \node[font=\bfseries,targetgreen!80!black] (n\r\c) at (\c*\w+\w/2, -\r*\h-\h/2) {\val};
      \else
        \node (n\r\c) at (\c*\w+\w/2, -\r*\h-\h/2) {\val};
      \fi
    }
  }

  \node[align=left,font=\footnotesize,draw=gray!50,fill=gray!5,rounded corners=3pt,inner sep=5pt] at (8, -2.5*\h) {
    \textbf{Finding Inverses:}\\
    To find $3^{-1}$, locate 3 in the left column,\\
    scan right to find the \textbf{1}, and look up:\\
    it sits in column 2, so $3^{-1} = 2$.\\
    Because $p=5$ is prime, every row\\
    is a permutation containing exactly one 1.
  };

  % Arrow for 3^-1 = 2 (curve cleanly below cell n31 so '3' is unobstructed)
  \draw[->,thick,quotientorange] (r3.south east) to[bend right=24] (n32.south west);
  \draw[->,thick,quotientorange] (n32.north) .. controls (\w*2.18, -\h*1.9) and (\w*2.18, -\h*1.1) .. (c2.south);
\end{tikzpicture}
\caption{The multiplication table for the nonzero elements of $\F_5$. The appearance of a 1 in every row and column confirms that every nonzero element has a unique multiplicative inverse. This structure defines a field.}
\label{ch05:fig-f5}
\end{figure}

Read off inverses by finding the $1$ in each row: $1^{-1}=1$, $2^{-1}=3$, $3^{-1}=2$,
$4^{-1}=4$. Every nonzero element has one, as a field demands. And the nonzero elements
$\{1,2,3,4\}$ form the multiplicative group $\Zns{5} = \F_5^{*}$; tracing powers of $2$
gives $2,4,3,1$, order $4$, so $2$ generates -- the multiplicative group of a finite field
is always cyclic, a fact we exploit constantly.
\end{example}

\begin{formalbox}
The cyclic structure of $\F_p^{*}$ is worth stating on its own because so much depends on
it: \emph{for every prime $p$, the multiplicative group $\F_p^{*}$ is cyclic of order
$p-1$.} A generator of it is called a \emph{primitive root} mod $p$. Its existence is what
guarantees Diffie--Hellman has a generator to stand on. We prove this carefully in
Chapter~6; for now, treat it as a reliable promise that prime moduli always
hand you a generator.
\end{formalbox}

\begin{focusbox}
Are there finite fields whose size is \emph{not} prime? Yes -- but only when the size is a
prime \emph{power} $p^k$. There is exactly one field of each prime-power size, written
$\F_{p^k}$, and none of any other size (no field of $6$ or $10$ elements exists). The
$p^k$ fields are not built as $\Zn{p^k}$ -- that ring has zero divisors -- but from
polynomials, which is the entire subject of the next chapter and the reason AES can treat
a byte as an element of $\F_{2^8}$.
\end{focusbox}

\section{A worked tour: the structures inside $\Zn{n}$}
\label{ch05:tour}

Let us pin all three structures onto a single modulus and watch them stack. Take $n=9$.

As a \textbf{ring}, $\Zn{9} = \{0,1,\dots,8\}$ with addition and multiplication mod $9$.
It is not a field: $9 = 3\cdot 3$ is a prime power but composite, and $3\cdot 3 = 9 \equiv
0$ makes $3$ a zero divisor.

Its \textbf{additive group} $(\Zn{9},+)$ has $9$ elements and is cyclic, generated by $1$
(or by any residue coprime to $9$, namely $1,2,4,5,7,8$).

Its \textbf{group of units} is $\Zns{9} = \{1,2,4,5,7,8\}$, the residues coprime to $9$, so
$\varphi(9)=6$. Is it cyclic? Trace powers of $2$ mod $9$:
\[
2,\;4,\;8,\;16\equiv 7,\;14\equiv 5,\;10\equiv 1.
\]
That is $2,4,8,7,5,1$ -- all six units, so $\ord(2)=6$ and $2$ is a generator. $\Zns{9}$ is
cyclic. (This matches the rule from Section~\ref{ch05:order}: $9 = 3^2$ is a prime power,
so its units form a cyclic group.) Contrast $\Zns{8}$ from
Example~\ref{ch05:ex-zns8}, which was \emph{not} cyclic -- the modulus matters.

\begin{walkbox}
Verify Lagrange inside $\Zns{9}$. Its size is $6$, so element orders must divide $6$, i.e.
lie in $\{1,2,3,6\}$. Check $\ord(8)$: $8 \equiv -1$, and $(-1)^2 = 1$, so $\ord(8)=2$.
Check $\ord(4)$: $4,16\equiv7,28\equiv1$, so $\ord(4)=3$. Both $2$ and $3$ divide $6$, as
required. The subgroup $\langle 4\rangle = \{1,4,7\}$ has size $3$, which divides $6$ --
two whole cosets tile the six units. Everything lines up with
Theorem~\ref{ch05:thm-lagrange}.
\end{walkbox}

\section{Toy code: groups by hand}
\label{ch05:code}

The fastest way to make these definitions concrete is to compute orders and hunt for
generators in code. The snippet below works in any $\Zns{n}$.

\begin{implementbox}
\textbf{Learning only -- never use this in real cryptography.} It is plain, variable-time,
and tiny on purpose. Real systems use vetted libraries with constant-time arithmetic and
carefully chosen, large parameters.
\begin{lstlisting}[language=Python]
from math import gcd

def units(n):
    "The group (Z/nZ)* : residues coprime to n."
    return [a for a in range(1, n) if gcd(a, n) == 1]

def order(a, n):
    "Smallest k > 0 with a^k = 1 mod n. Assumes gcd(a, n) = 1."
    k, x = 1, a % n
    while x != 1:
        x = (x * a) % n
        k += 1
    return k

def generators(n):
    "Elements whose order equals |(Z/nZ)*| = phi(n); empty if non-cyclic."
    U = units(n)
    phi = len(U)
    return [a for a in U if order(a, n) == phi]

# Check Lagrange: every order divides phi(n).
n = 9
U = units(n)                       # [1, 2, 4, 5, 7, 8]
phi = len(U)                       # 6
orders = {a: order(a, n) for a in U}
assert all(phi % orders[a] == 0 for a in U)   # Lagrange holds
print("phi =", phi)
print("orders =", orders)          # {1:1, 2:6, 4:3, 5:6, 7:3, 8:2}
print("generators =", generators(n))          # [2, 5]

# A non-cyclic group has no generator at all:
print("generators of 8:", generators(8))      # []  (Z/8Z)* is not cyclic
\end{lstlisting}
The assertion is Lagrange's theorem turned into a runtime check: no matter which $n$ you
plug in, every element's order divides $\varphi(n)$. The empty generator list for $n=8$
is Example~\ref{ch05:ex-zns8} reappearing as data.
\end{implementbox}

\section{A pitfall: small subgroups and where Lagrange bites back}
\label{ch05:pitfall}

Lagrange is usually your friend, but an attacker can turn the same theorem against a
careless implementation. The danger is the \emph{small subgroup}.

\begin{attackbox}
\textbf{Small-subgroup confinement.} Diffie--Hellman in $\Zns{p}$ is supposed to live in a
large prime-order subgroup. But $\Zns{p}$ has size $p-1$, which is even, so by Lagrange it
\emph{always} contains a subgroup of order $2$ -- namely $\{1, p-1\}$, since
$(p-1)^2 \equiv 1$. More generally, every divisor of $p-1$ gives a subgroup of that size.

Now suppose Alice computes a shared secret as $K = Y^a \bmod p$, where $Y$ is a value Bob
sent. If Alice does not validate $Y$, Eve (impersonating Bob) can send a $Y$ of small
order -- say order $2$, so $Y \in \{1, p-1\}$. Then $K = Y^a$ can take only $2$ possible
values regardless of Alice's secret $a$, and which one it is leaks the parity of $a$
(Figure~\ref{ch05:fig-smallsub}).
Repeat with subgroups of order $3, 5, \dots$ and Eve peels off $a$ modulo each small
factor, then reassembles $a$ by the Chinese Remainder Theorem. The exponent's secrecy
collapses one small subgroup at a time.

\textbf{The fix is pure Lagrange.} Choose $p = 2q+1$ with $q$ prime (a \emph{safe
prime}), so the only subgroup sizes are $1,2,q,2q$ -- no small factors to confine into --
and work in the order-$q$ subgroup. Then \emph{validate} every received element: reject
anything not in that subgroup, e.g. by checking $Y^q \equiv 1$ and $Y \ne 1$. Elliptic-curve
deployments do the analogous ``point validation'' and choose curves of nearly prime order
(small \emph{cofactor}) for exactly the same reason.
\end{attackbox}

\begin{applicationbox}
This is not a textbook hypothetical. The order-$q$ subgroup of a safe-prime $\Zns{p}$ is
the setting of classic Diffie--Hellman as standardized for TLS and IPsec. RFC~7919 ships
fixed, vetted safe-prime groups (``ffdhe2048'', ``ffdhe3072'', \dots) precisely so
implementers stop rolling their own parameters and stop tripping over small subgroups.
Modern TLS 1.3 leans instead on elliptic-curve groups such as Curve25519, whose designer
Daniel Bernstein chose the curve so that the cofactor is a tiny $8$ and the
order-validation problem largely takes care of itself. In every case the design rule is
the same sentence: \emph{know the order of your group, and make sure every element you
accept lives in the big prime-order piece of it.}
\end{applicationbox}

\begin{figure}[H]
\centering
\begin{tikzpicture}[font=\small,>=stealth]
  % Party Nodes
  \node[draw=aliceblue!80,fill=aliceblue!15,thick,rounded corners=3pt,minimum width=2.6cm,minimum height=0.7cm,font=\bfseries] (alice) at (0,0) {Alice};
  \node[draw=evered!80,fill=evered!15,thick,rounded corners=3pt,minimum width=2.6cm,minimum height=0.7cm,font=\bfseries] (eve) at (8.5,0) {Eve (Active Impostor)};

  % Timelines
  \draw[dashed,aliceblue!70,thick] (alice.south) -- (0,-3.4);
  \draw[dashed,evered!70,thick] (eve.south) -- (8.5,-3.4);

  % Step 1: Eve injects Y = p - 1
  \draw[->,very thick,evered] (8.5,-0.7) -- node[above,font=\scriptsize\bfseries]{(1) Injects $Y = p-1 \equiv -1 \pmod p$ (order $2$)} (0,-0.7);

  % Step 2: Alice collapses key
  \node[draw=aliceblue!70,fill=white,rounded corners=2pt,align=left,font=\tiny,inner sep=4pt] at (0,-1.65) {
    (2) Computes shared key:\\
    $K = Y^a \equiv (-1)^a \pmod p$\\[2pt]
    $\to K \equiv 1 \pmod p$ (if $a$ even)\\
    $\to K \equiv -1 \pmod p$ (if $a$ odd)
  };

  % Step 3: Leakage back to Eve
  \draw[->,thick,dashed,evered] (0,-2.6) -- node[below,font=\tiny\bfseries]{(3) Encrypted response leaks $a \bmod 2$} (8.5,-2.6);

  % Defense Box
  \node[draw=targetgreen!70,fill=targetgreen!8,rounded corners=3pt,thick,align=left,font=\tiny,inner sep=5pt] at (4.25,-3.8) {
    \textbf{The Defense (Lagrange in Action):}\\
    Alice checks subgroup membership before computing $K$: test $Y^q \equiv 1 \pmod p$\\
    and $Y \not\equiv 1$. Order-2 elements are instantly rejected, closing the confinement attack.
  };
\end{tikzpicture}
\caption{Small-subgroup confinement attack on unvalidated Diffie-Hellman. Eve sends an element $Y = p-1$ of order 2. Alice's exponentiation collapses $K$ into a 1-bit value revealing $a \bmod 2$. Group membership validation blocks the attack by checking that received elements live strictly in the large prime-order subgroup.}
\label{ch05:fig-smallsub}
\end{figure}

\section{How the three structures power the rest of the book}
\label{ch05:roadmap}

Step back and see the layering, because the rest of the book is built on it.

\textbf{Groups} give you the discrete-log world: a cyclic group with a hard-to-invert
exponentiation map. That single idea drives Diffie--Hellman key exchange
(Chapter~18), ElGamal and DSA (Chapter~20), and elliptic-curve cryptography
(Chapter~21), where the group is points on a curve instead of residues but the algebra is
the same.

\textbf{Rings} give you $\Zn{n}$ with both operations, which is the natural home of RSA
(Chapter~19): encryption is $m \mapsto m^e \bmod n$ for a \emph{composite} $n=pq$, and the
correctness proof is Euler's theorem -- our Corollary~\ref{ch05:cor-euler-fermat} -- applied
in the ring $\Zn{n}$. The composite modulus, with its zero divisors, is the whole point:
factoring $n$ is what is hard.

\textbf{Fields} give you exact, reversible arithmetic with no zero divisors. Prime fields
$\F_p$ underlie elliptic curves (the curve lives over a field). Prime-power fields
$\F_{2^k}$ underlie the byte-level mixing in AES and the carry-less arithmetic in
GCM authentication. The next chapter builds those fields from the ring of polynomials --
yet another instance of ``take a ring, quotient by something irreducible, and a field
pops out,'' the exact same move that turned $\Z$ into $\F_p$ by quotienting by a prime.

\begin{intuitionbox}
If you remember one sentence from this chapter, make it this: \emph{a group lets you
combine and undo, a ring adds a second operation, and a field lets you divide; crypto
picks whichever structure makes the easy direction cheap and the hard direction
hopeless.} Everything else -- generators, orders, Lagrange, zero divisors -- is bookkeeping
in service of that choice.
\end{intuitionbox}

\section{Exercises}
\label{ch05:exercises}

\begin{exercise}[Group or not?]
\label{ch05:ex-group-or-not}
For each, decide whether it is a group; if not, name the failing axiom.
(a) The even integers under addition.
(b) The odd integers under addition.
(c) $\{1,2,3,4\}$ under multiplication mod $5$.
(d) $\{0,1,2,3\}$ under multiplication mod $4$.
(e) The set $\Zns{n}$ under multiplication mod $n$, for general $n$.
\end{exercise}

\begin{exercise}[Orders by hand]
\label{ch05:ex-orders-hand}
Work entirely mod $11$. (a) Compute the order of every element of $\Zns{11}$. (b) List all
generators. (c) Confirm that each order you found divides $10$, and explain which theorem
guarantees this in advance.
\end{exercise}

\begin{exercise}[Cyclic or not]
\label{ch05:ex-cyclic-or-not}
Determine whether $\Zns{n}$ is cyclic for $n = 7, 8, 9, 10, 12, 16$. For each cyclic case
exhibit a generator; for each non-cyclic case, show that every element's order is strictly
less than $\varphi(n)$ and identify the largest order that occurs.
\end{exercise}

\begin{exercise}[Subgroups and Lagrange]
\label{ch05:ex-subgroups}
In $\Zns{13}$ (which has $12$ elements): (a) Find the cyclic subgroup generated by $3$ and
state its size. (b) Find a subgroup of size $4$ and one of size $6$. (c) Explain why no
subgroup of size $5$ can exist, citing Lagrange.
\end{exercise}

\begin{exercise}[Euler from Lagrange]
\label{ch05:ex-euler-proof}
Using only Corollary~\ref{ch05:cor-glag} (``$g^{\abs{G}} = e$''), prove Euler's theorem:
if $\gcd(a,n)=1$ then $a^{\varphi(n)} \equiv 1 \pmod n$. Then specialize to deduce Fermat's
little theorem. State precisely which group you are applying the corollary to and why $a$
belongs to it.
\end{exercise}

\begin{exercise}[Zero divisors and cancellation]
\label{ch05:ex-zerodiv}
(a) List all zero divisors in $\Zn{18}$. (b) Show $\gcd(a,18) > 1$ exactly characterizes
them. (c) Exhibit a concrete failure of cancellation: find $a \ne 0$ and $x \ne y$ in
$\Zn{18}$ with $a x \equiv a y$. (d) Explain why the same cannot happen in $\F_{19}$.
\end{exercise}

\begin{exercise}[Fields and non-fields]
\label{ch05:ex-fields}
(a) Prove that $\Zn{n}$ is a field iff $n$ is prime, filling in every step of
Theorem~\ref{ch05:thm-prime-field}. (b) Build the full multiplication table of $\F_7$ and
read off the inverse of each nonzero element. (c) Explain in one sentence why no field of
order $6$ exists.
\end{exercise}

\begin{exercise}[Break a careless Diffie--Hellman]
\label{ch05:ex-small-subgroup}
Let $p = 23$, so $p - 1 = 22 = 2 \cdot 11$. (a) Find an element of order $2$ and an
element of order $11$ in $\Zns{23}$. (b) Suppose Alice computes $K = Y^a \bmod 23$ without
validating $Y$, and Eve sends the order-$2$ element as $Y$. Show that $K$ reveals
$a \bmod 2$. (c) State the safe-prime fix and the validation check that defeats this
attack, connecting it to the small-subgroup discussion in
Section~\ref{ch05:pitfall}.
\end{exercise}

\begin{exercise}[Generators of the additive group]
\label{ch05:ex-additive-gen}
Prove that $k$ generates the additive group $\Zn{n}$ if and only if $\gcd(k,n)=1$. (Hint:
the multiples of $k$ are $\{0, k, 2k, \dots\} \bmod n$; relate the number of distinct
values to $\ord$ of $k$ and to $\gcd(k,n)$.)
\end{exercise}

\begin{exercise}[Count the generators]
\label{ch05:ex-count-gen}
A cyclic group of order $n$ has exactly $\varphi(n)$ generators. (a) Verify this for
$\Zns{7}$ ($n=6$) by listing its generators and counting. (b) Verify for $\Zns{9}$
($n=6$). (c) Use the formula to predict, without enumeration, how many primitive roots
$\Zns{p}$ has when $p = 2q+1$ is a safe prime; evaluate it for $p=23$.
\end{exercise}
